\documentclass[11pt]{article}
\usepackage[margin=1in]{geometry}
\usepackage[T1]{fontenc}
\usepackage{lmodern}
\usepackage{amsmath,amssymb,amsthm,mathtools}
\usepackage{algorithm}
\usepackage[noend]{algpseudocode}
\usepackage{listings}
\usepackage{booktabs}
\usepackage{array}
\usepackage{ebproof}
\usepackage{tabularx}
\usepackage{longtable}
\newcolumntype{L}[1]{>{\raggedright\arraybackslash}p{#1}}
\usepackage{enumitem}
\usepackage{xcolor}
\usepackage{flafter}
\usepackage{tikz}
\usetikzlibrary{arrows.meta,positioning,patterns,shapes.geometric,fit,calc}
\usepackage[colorlinks=true,linkcolor=blue!50!black,citecolor=blue!50!black,urlcolor=blue!50!black]{hyperref}
\usepackage[capitalise,nameinlink]{cleveref}

\lstset{basicstyle=\ttfamily\small,columns=fullflexible,keepspaces=true,
        frame=single,framerule=0.3pt,rulecolor=\color{black!30},
        xleftmargin=1em,xrightmargin=1em,breaklines=true}

\emergencystretch=3em
\theoremstyle{plain}
\newtheorem{theorem}{Theorem}
\newtheorem{proposition}[theorem]{Proposition}
\newtheorem{lemma}[theorem]{Lemma}
\newtheorem{corollary}[theorem]{Corollary}
\theoremstyle{definition}
\newtheorem{definition}[theorem]{Definition}
\newtheorem{example}[theorem]{Example}
\newtheorem{observation}[theorem]{Observation}
\theoremstyle{remark}
\newtheorem{remark}[theorem]{Remark}
\crefname{observation}{Observation}{Observations}

\begin{document}
\section{Proof-based Interpolation}
We perform proof-based interpolation in the style of McMillan, ``Interpolants from Z3 Proofs.'' We may view this method as a set of Z3 proof-directed simplification rules to simplify the interpolation goal as much as possible, before delegating interpolant computation for a simpler formula into a secondary interpolating prover. Recall that for a given pair of formulas $(A, B)$ a \textbf{local} sequent is a sequent entirely in $\mathcal{L}(A)$ or in $\mathcal{L}(B)$. A collection of formulas $\Gamma$ is called \textbf{strict} if each formula of $\Gamma$ is local but $\Gamma$ itself features both formulas referencing $A$'s vocabulary and formulas referencing $B$'s vocabulary.  
We assume the proof calculus for Z3 can be normalized to having the following form: 

\begin{center}
\vspace*{0.1in}
\begin{prooftree}
    \hypo{\Gamma, \Delta \text{ are local}}
    \Infer1[LOCAL]{\Gamma \vdash \Delta}
\end{prooftree}

\vspace*{0.1in}
\begin{prooftree}
    \hypo{\Gamma \vdash \Delta, p}
    \hypo{\Gamma' \vdash \Delta', \neg p}
    \Infer2[RES]{\Gamma,\Gamma'\vdash \Delta, \Delta'}
\end{prooftree}

\vspace*{0.1in}
\begin{prooftree}
    \hypo{\,\,\Gamma, \Gamma' \vdash}
    \Infer1[CONTRA]{\Gamma' \vdash \neg \Gamma'}
\end{prooftree}

\vspace*{0.1in}
\begin{prooftree}
    \hypo{\Gamma,\Delta \text{ are strict}}
    \Infer1[LEMMA]{\Gamma\vdash \Delta}
\end{prooftree}
\end{center}

Note that the LEMMA rule is essentially a relaxed version of the LOCAL rule, and the LEMMA leafs are exactly the proof obligations in need of delegating to a secondary interpolation prover.

\paragraph{The interpolating calculus.}
An \textbf{interpolation} is a sequent of form $(A, B) \vdash \Delta \,[\phi]$. Let $\Delta_A$ denote the parts of $\Delta$ that occur local in $A$, and let $\Delta_B$ denote the parts of $\Delta$ that occur local in $B$. We say this sequent is valid when
\begin{enumerate}
    \item $A \wedge \neg \Delta_A \vDash \phi$
    \item $B \wedge \neg \Delta_B \wedge \phi$ is UNSAT
    \item $\mathcal{L}(\phi)\subset \mathcal{L}(A)\cap\mathcal{L}(B)$
\end{enumerate}
We may design sequent interpolants via the following set of rules
\begin{gather*}
% LocalA
\begin{prooftree}
  \infer[
    left label=\textsc{LocalA},
    right label={$\bigwedge A \models_{\mathcal T} \bigvee\Delta,
                  \quad \Delta \subseteq \mathcal L(A)$}
  ]0{
    (A,B) \vdash \Delta \; [\bigvee\Delta_B]
  }
\end{prooftree}
\\[1.5ex]
% LocalB
\begin{prooftree}
  \infer[
    left label=\textsc{LocalB},
    right label={$\bigwedge B \models_{\mathcal T} \bigvee\Delta,
                  \quad \Delta \subseteq \mathcal L(B)$}
  ]0{
    (A,B) \vdash \Delta \; [\text{\textsc{True}}]
  }
\end{prooftree}
\\[1.5ex]
% ResA
\begin{prooftree}
  \hypo{(A,B) \vdash \Delta,p \; [\phi]}
  \hypo{(A,B) \vdash \Delta',\neg p \; [\phi']}
  \infer[
    left label=\textsc{ResA},
    right label={$p \in \mathcal L(A) \setminus \mathcal L(B)$}
  ]2{
    (A,B) \vdash \Delta \cup \Delta' \; [\phi \lor \phi']
  }
\end{prooftree}
\\[1.5ex]
% ResB
\begin{prooftree}
  \hypo{(A,B) \vdash \Delta,p \; [\phi]}
  \hypo{(A,B) \vdash \Delta',\neg p \; [\phi']}
  \infer[
    left label=\textsc{ResB},
    right label={$p \in \mathcal L(B)$}
  ]2{
    (A,B) \vdash \Delta \cup \Delta' \; [\phi \land \phi']
  }
\end{prooftree}
\\[1.5ex]
% Contra
\begin{prooftree}
  \hypo{(A,B) \vdash \Delta \; [\phi]}
  \infer[
    left label=\textsc{Contra}
  ]1{
    (A,B) \vdash \Delta \; [\phi]
  }
\end{prooftree}
\end{gather*}

In what follows, we will show that as long as a given theory $Th$ supports model-based projection, we can design a model-based projection loop that handles all the \textsc{Lemma} nodes encountered in the above interpolating calculus.

\section{Model-based Interpolation}

This section introduces a model-based interpolation procedure that can be used on its own to generate interpolants for pairs of formulas, or be combined as a leaf-level solver for the LEMMA nodes inside the proof-based rules above.

Recall the definition of Craig interpolation:
\begin{definition}[Craig Interpolants]

Given a formula $F(X, Z)$ and a formula $G(Z, Y)$ such that $F \wedge G$ is UNSAT, a \textbf{Craig interpolant} $I$ is a formula over vocabulary $Z$ such that (i) $F \vDash I$; (ii) $I \wedge G$ is UNSAT.
\end{definition}


%We assume that if $F$ is an UNSAT formula, then we can extract an implicant of a cube $F$ called the \textbf{UNSAT core} of $F$. An UNSAT core $C$ for $F$ is a subset of the conjuncts in $F$ such that $C$ is still UNSAT. 

\begin{definition}[Model-based Projection]
    Let $F$ be a conjunction of terms over vocabulary $X$, and let $M$ be a model of $F$. A \textbf{model-based projection} of $F$ on $S\subset X$ using $M$ is a cube $P$ such that
    (i) $P$ is over vocabulary $X\setminus S$; (ii) $M\vDash P$; (iii) $P \vDash \exists S.F$. We denote the formula $P$ by $Mbp(F,S,M)$. A model-based projection operator $Mbp(-)$ has the \textbf{finite-image property} if there are only finitely many model-based projection formulas for a given $F$ and $M$.
\end{definition}

We asume for a formula $F$, the following primitives are available:
\begin{enumerate}
    \item If $F$ is SAT, then $Mbp(F,S,M)$ is the model-based projection of $\exists S. F$ at $M$;
    \item If $F$ is UNSAT, then $Core(F)$ is the minimal UNSAT core of $F$.
\end{enumerate}

We claim the following algorithm using the above primitives produces a Craig interpolant, should it terminate:

\begin{algorithmic}
    \Procedure{MBI}{$A$, $B$}
        \State $I \gets [];$
        \State $F\gets A;$
        \State $S \gets \mathcal{L}(A)\setminus \mathcal{L}(B);$
        \While{$\text{Check}(F) \neq \text{UNSAT}$}
            \State $M \gets \text{GetModel}(F);$
            \State $P \gets Mbp(A,M,S);$
            \If{$\text{Check}(B \wedge P)=\text{SAT}$}
                \Return $(A, B)$ is SAT;
            \Else 
                \State $C \gets Core(B \wedge P);$
                \State $F \gets F \wedge \neg (\bigwedge C);$
                \State $I \gets I \cup C;$
                \EndIf
        \EndWhile 
        \Return $\bigvee_{C \in I} (\bigwedge C)$;
    \EndProcedure 
\end{algorithmic}

It is critical to note that any model $M$ of $F$ is also a model $M$ of $A$, since $F$ is always a strengthening of $A$. We have the following theorem:

\begin{theorem}
    If the MBI($A, B$) procedure returns a formula $I$, then $I$ is a Craig interpolant of $(A, B)$.
\end{theorem}

\begin{proof}
Observe the following facts:
\begin{enumerate}
    \item At every iteration of MBI's inner loop, $A \vDash C$;
    \item At every iteration of MBI's inner loop, $C \wedge B$ is UNSAT.
\end{enumerate}

    If MBI($A, B$) returns such a formula $I$, then its inner loop must have terminated at some step $k$, with
    \[
    F = A \wedge \neg C_1 \wedge \dots \wedge\neg C_n
    \]
    Since $A \vDash A$ and $A \vDash C_i$, we conclude that $A \vDash \bigvee_{i=1}^n C_i$. Since $B$ and each $C_i$ is UNSAT, $B \wedge \bigvee_{i=1}^nC_i = \bigvee_{i=1}^n B \wedge C_i$ is UNSAT. This concludes the proof.
\end{proof}

Furthermore, the above procedure terminates if the model-based projection operator has the finite-image property.

The above algorithm shows up in various places, including in the paper ``Model Based Interpolation for Uninterpreted
Functions and Integer Linear Arithmetic'' by Nikolaj Bjorner, Arie Gurfinkel, Sharon Shoham and Yakir Vizel, LPAR 2023.


\end{document}

